\documentclass[11pt]{article}

\usepackage[margin=1in]{geometry}
\usepackage{graphicx}
\usepackage{booktabs}
\usepackage{amsmath}
\usepackage{amssymb}
\usepackage{caption}
\usepackage{subcaption}
\usepackage[colorlinks=true,linkcolor=blue,citecolor=blue,urlcolor=blue]{hyperref}

\title{Latent Dynamics of a Small MNIST Classifier: Geometry and Topology over Training}
\author{Generated by Claude (sonnet)}
\date{2026-10-14}

\begin{document}
\maketitle

\begin{abstract}
We study how the latent representation of a small MLP classifier on MNIST evolves
during training, tracking a fixed probe set at 50-step intervals and computing
geometric (centroid separation, Fisher-like discriminant ratio) and topological
(persistent-homology, magnitude-function) summaries at every snapshot. Three
experiments vary (i) latent dimension $D \in \{2,4,8,16,32\}$, (ii) training
objective (supervised probe, autoencoder reconstruction, or both), and (iii)
a latent normalization layer (none, $\ell_2$-projection to the unit sphere, or
BatchNorm), each at 3 seeds. We find that test accuracy alone is a poor summary
of what is happening in the latent space: the raw latent norm inflates roughly
20--40$\times$ over training and must be factored out before geometric metrics
are meaningful; once this scale is removed, accuracy-matched configurations
turn out to have substantially different latent geometry and topology.
Supervision collapses the latent into topologically simple, well-separated
class clusters, while reconstruction yields a diffuse, topologically rich
manifold with much weaker class alignment. Capacity has essentially no effect
on accuracy above $D{=}4$, and test accuracy converges almost immediately
(within $\sim$4--6\% of training steps for $D\ge4$); the latent's
\emph{geometry}, by contrast, keeps reorganizing for much longer
($\sim$22--36\% of training before the Fisher ratio reaches 90\% of its
final value), and the normalized Fisher ratio is itself non-monotonic in $D$
(peaking near $D{=}4$) while persistent-homology features (H1 entropy, H1
lifetime sum) increase monotonically with $D$ up to $D{=}16$ and then
saturate -- $D{=}32$ adds essentially no further geometric or topological
structure over $D{=}16$. Constraining the latent to the unit sphere
($\ell_2$ normalization) achieves the same accuracy as the unconstrained
probe at essentially zero cost while producing far stronger angular class
separation (Fisher ratio $\approx$15 vs.\ $\approx$3.6) and a much larger
scale-averaged magnitude (PMag $\approx$53 vs.\ $\approx$8); since the
Fisher ratio is mathematically invariant to isotropic rescaling, this gain
cannot be an artifact of the latent's raw scale and instead reflects genuine
representational reshaping into tighter, better angularly-separated class
clusters, induced by training the encoder under an anisotropic, per-sample
normalization constraint.
\end{abstract}

\section{Introduction and Setup}

Test accuracy and training loss describe how well a network solves its task,
but they say little about the internal representation that solves it. Two
models with identical accuracy can organize their latent space very
differently -- more or less compactly, more or less aligned to class labels,
with simpler or richer topology. This report tracks those internal
properties directly, over the course of training, for a small MNIST
classifier, and asks three questions: (1) how does the latent space change
shape as training proceeds, separately from how it changes scale; (2) how
does latent capacity (dimension $D$) affect the geometry and topology of the
learned representation, independently of its effect on accuracy; and (3) how
does the training objective (supervised vs.\ reconstructive) and an explicit
normalization constraint on the latent change that representation.

\paragraph{Model.} The model is a small MLP encoder, $784 \to 256 \to 64 \to
D$, mapping a flattened MNIST image to a $D$-dimensional latent vector $z$.
An optional normalization layer is applied to $z$ (identity / $\ell_2$
projection onto the unit sphere / BatchNorm). A linear probe head (readout)
maps $z$ to 10 class logits.

\paragraph{Objectives.} Three training objectives are compared, all on the
same encoder architecture:
\begin{itemize}
  \item \textbf{probe}: cross-entropy loss on the probe head's logits;
        gradients flow through the probe head into the encoder, so the
        encoder is shaped directly by the classification objective.
  \item \textbf{autoencoder}: a decoder reconstructs the input from $z$ and
        the encoder is trained only to minimize reconstruction MSE; the
        probe head is still trained (for monitoring), but it reads a
        \emph{detached} copy of the latent, so its accuracy measures the
        linear separability of a representation that was never shaped by
        the classification signal.
  \item \textbf{probe\_recon}: the sum of the cross-entropy and
        reconstruction losses, so both signals shape the encoder jointly.
\end{itemize}

\paragraph{Probe-set snapshotting.} A fixed, class-balanced set of $N{=}1000$
held-out test samples is chosen once per run and never changes. Every 50
optimizer steps, the current encoder is used to compute the latent vectors
for exactly this probe set, and the full battery of geometric/topological
metrics (Section~\ref{sec:methods}) is computed on that snapshot. Because the
sample set and order are fixed, snapshots are directly comparable across
steps and across runs. Each run trains for 5 epochs ($\approx$2345 optimizer
steps at batch size 128), giving 48 snapshots per run.

\paragraph{Experiments.} All experiments share the snapshot/metric pipeline
above.
\begin{enumerate}
  \item \textbf{Baseline}: a single $D{=}16$, objective=probe,
        normalization=none run, examined in detail to establish the
        scale-vs-shape distinction and the overall shape of the training
        trajectory.
  \item \textbf{Capacity sweep}: $D \in \{2,4,8,16,32\}$, objective=probe,
        normalization=none, $\times$ 3 seeds (15 runs).
  \item \textbf{Objective contrast}: $D{=}16$, normalization=none, three
        objectives (probe, autoencoder, probe\_recon) $\times$ 3 seeds
        (9 runs).
  \item \textbf{Normalization contrast}: $D{=}16$, objective=probe, three
        normalization modes (none, $\ell_2$, BatchNorm) $\times$ 3 seeds
        (9 runs).
\end{enumerate}
In total 27 runs (reported in Table~\ref{tab:results}); error bands/bars
throughout are $\pm 1$ standard deviation across the 3 seeds.

\section{Methods}
\label{sec:methods}

\subsection{Scale vs.\ shape: normalization}

Over training, the raw latent space inflates roughly isotropically (its mean
norm grows by more than an order of magnitude, Section~\ref{sec:baseline}).
Since every distance-based or topological metric below is a function of
pairwise Euclidean distances, an isotropic scale change multiplies all of
them by (approximately) the same factor and would otherwise dominate and
mask genuine reorganization of the latent \emph{shape} (how class clusters
move relative to one another). To separate the two, every
distance/topology metric is computed on a per-snapshot normalized point
cloud,
\[
  z_i' = \frac{z_i}{\bar{d}}, \qquad
  \bar d = \operatorname*{mean}_{i<j} \lVert z_i - z_j \rVert ,
\]
i.e.\ every coordinate of every point is divided by the mean pairwise
Euclidean distance of the cloud at that snapshot. This is an isotropic
rescaling by a single scalar, so it preserves all angles and all relative
distance ratios (hence the full shape of the point cloud, including its
persistence diagrams up to a uniform rescaling of the filtration parameter)
while removing the absolute scale. The raw (un-normalized) scale is tracked
separately via \texttt{raw\_mean\_norm}, the mean of $\lVert z_i \rVert$ over
the probe set.

Not every metric needs this treatment equally: distances, persistence
diagrams, and the magnitude function are all scale-\emph{sensitive} by
construction (an isotropic rescaling $z \mapsto cz$ rescales every pairwise
distance by $c$, hence rescales lifetimes, lifetime sums, and the magnitude
curve), so normalizing before computing them is essential. The Fisher-like
separation ratio below is the one exception: it is a ratio of sums of
squared distances ($S_{\text{between}}/S_{\text{within}}$), and both
numerator and denominator scale by $c^2$ under $z \mapsto cz$, so the ratio
is exactly scale-\emph{invariant} (verified empirically on every snapshot of
every run in this study: fisher\_ratio computed on raw latents agrees with
fisher\_ratio computed on normalized latents to within $\sim 10^{-15}$
relative error). We still report it on normalized latents for consistency
with the other metrics, but this invariance is important for interpreting
Section~\ref{sec:norm}: any Fisher-ratio change under a normalization layer
cannot be explained by a change in raw scale alone.

\subsection{Class geometry}

Let $z_1,\dots,z_N \in \mathbb{R}^D$ be the (normalized, unless noted) probe
latents with labels $y_i \in \{0,\dots,9\}$, and let $\mu_c =
\operatorname{mean}_{i: y_i = c} z_i$ be the centroid of class $c$ and
$\mu = \operatorname{mean}_i z_i$ the global centroid.

\paragraph{Separation.} \texttt{separation\_mean} is the mean, over all
class pairs $a<b$, of the inter-centroid distance $\lVert \mu_a - \mu_b
\rVert$.

\paragraph{Fisher-like separation ratio.} Using the classical
multivariate (trace-of-scatter) generalization of the one-dimensional Fisher
discriminant ratio,
\[
  S_{\text{between}} = \sum_c n_c \lVert \mu_c - \mu \rVert^2 , \qquad
  S_{\text{within}} = \sum_c \sum_{i: y_i = c} \lVert z_i - \mu_c \rVert^2 ,
  \qquad
  \text{fisher\_ratio} = \frac{S_{\text{between}}}{S_{\text{within}}} ,
\]
with $n_c$ the number of probe points in class $c$. This needs no
eigendecomposition (unlike LDA's generalized eigenvalue criterion) and is
larger when classes are better separated relative to their own spread.

\subsection{Topological metrics (persistent homology)}

Persistence diagrams are computed with \texttt{ripser} on the Vietoris--Rips
filtration of the (normalized) probe point cloud, up to homological degree 1
(connected components H0 and independent loops H1). For a diagram with bars
$(b_i,d_i)$ (birth, death), we use two summaries:

\paragraph{Alpha-weighted lifetime sum.}
\[
  L_\alpha(\text{dgm}) = \sum_i (d_i - b_i)^\alpha ,
\]
summed over all (capped-)finite bars. $\alpha{=}1$ is exactly the \emph{total
persistence} (sum of lifetimes); $\alpha{=}2$ emphasizes long-lived
features more strongly, and $\alpha{=}0.5$ dampens the influence of a few
very long-lived bars. We report $L_1$ for H1 (\texttt{h1\_lifetime\_sum\_alpha1.0})
as the primary topological-richness summary. The single infinite-death H0
bar (the component that never merges) is capped at the largest finite death
value in the diagram before being included, so that H0 totals remain finite
and comparable across snapshots instead of discarding that bar.

\paragraph{Persistence entropy.} Following the standard normalized-Shannon-entropy
definition of persistence entropy (Rucco et al.; Atienza et al., ``Stability
properties of persistent entropy''), with $l_i = d_i - b_i$ the lifetime of
bar $i$ and $L = \sum_i l_i$,
\[
  p_i = \frac{l_i}{L}, \qquad
  \mathrm{PE}(\text{dgm}) = -\sum_i p_i \log p_i .
\]
Low entropy means persistence is concentrated in a few dominant bars
(simple, few-feature topology); high entropy means persistence is spread
across many comparably long-lived bars (rich topology).

\subsection{Magnitude function and PMag}

The \emph{magnitude} of a finite metric space $(X,d)$ (Leinster, 2013,
``The magnitude of metric spaces,'' Documenta Mathematica 18) is defined via
the similarity matrix $Z_{ij} = \exp(-t\, d(x_i,x_j))$ at scale $t>0$: if $Z$
is invertible, solving $Zw = \mathbf{1}$ for the weighting $w$ gives
\[
  \mathrm{Mag}(t) = \sum_i w_i = \sum_{i,j} (Z^{-1})_{ij} .
\]
As $t \to \infty$, $\mathrm{Mag}(t) \to |X|$ (every point resolved as
distinct); as $t \to 0$, nearby points become redundant and contribute less.
The magnitude \emph{function} $t \mapsto \mathrm{Mag}(t)$ is thus a
multi-scale summary of the point cloud's ``effective size,'' directly
analogous to a persistence diagram's multi-scale summary of topology.
Govc \& Hepworth (``Persistent magnitude,'' J. Pure Appl. Algebra 225 (2021);
arXiv:2010.09776) make this analogy precise, defining persistent magnitude
of a filtered Rips complex as a magnitude-valued invariant that plays the
same role across a filtration that persistence diagrams play for homology.

We implement $\mathrm{Mag}(t)$ by solving $Zw=\mathbf 1$ (robustly: a tiny
diagonal jitter, falling back to least-squares and then to the Moore--Penrose
pseudo-inverse for ill-conditioned $Z$) at $n_\text{scales}{=}12$ scales
$t$ drawn from a data-driven range $[t_{\min}, t_{\max}] = [1/d_{95},
1/d_{5}]$, where $d_5, d_{95}$ are the 5th/95th percentiles of the nonzero
pairwise-distance distribution of the cloud -- keeping the scale range
comparable across snapshots even as the raw latent scale drifts. \textbf{PMag} is
then the scale-averaged magnitude, i.e.\ the trapezoidal integral of
$\mathrm{Mag}(t)$ normalized by the scale range width,
\[
  \mathrm{PMag} = \frac{1}{t_{\max}-t_{\min}} \int_{t_{\min}}^{t_{\max}}
  \mathrm{Mag}(t)\,dt ,
\]
a single scalar summarizing the whole magnitude curve (rather than an
arbitrary single-scale sample), in the same ``area under a scale-indexed
curve'' spirit as the alpha-weighted lifetime sums above.

\section{Results}

\subsection{Baseline: scale inflation and the shape ``elbow''}
\label{sec:baseline}

In the baseline run ($D{=}16$, objective=probe, normalization=none), the raw
mean latent norm grows from $0.57$ at initialization to $23.7$ at the final
step -- a roughly 40$\times$ isotropic inflation, with most of the growth
happening abruptly within the first logging interval (step 50: already
$16.2$) and a slower continued drift afterward (step 2345: $23.7$). Because
every raw distance-based metric scales with this norm, \texttt{raw\_separation\_mean}
tracks \texttt{raw\_mean\_norm} almost exactly in shape
(Figure~\ref{fig:scale-shape}, left panel) and is therefore dominated by
scale, not by any change in how classes are arranged relative to each other.

The scale-robust, normalized metrics tell a different and more informative
story (Figure~\ref{fig:scale-shape}, right panel). The normalized Fisher
ratio rises sharply from $0.19$ at step 0 to $\approx 3.1$--$3.4$ by step
300--500 (roughly the first epoch), then continues to refine only slowly,
reaching $3.72$ at the final step (2345). Normalized \texttt{separation\_mean}
shows the same pattern: a sharp rise to $\approx0.90$--$0.95$ within the
first $\sim$500 steps, then a much slower approach to its final value of
$0.964$. In other words, the real geometric reorganization of the latent
space -- not just its overall growth -- happens as a sharp ``elbow'' within
roughly the first epoch, followed by slow refinement over the remaining four
epochs. Final-step test accuracy for this run is $0.975$.

\begin{figure}[htbp]
  \centering
  \includegraphics[width=0.95\textwidth]{figures/baseline_scale_vs_shape.png}
  \caption{Baseline run ($D{=}16$, probe, no normalization). Left: raw
  (un-normalized) latent-space inflation -- mean norm and mean inter-centroid
  separation both grow $\sim$isotropically over training. Right: the same
  run's \emph{normalized} separation\_mean and fisher\_ratio, which isolate
  shape from scale and reveal a sharp early-training elbow (first
  $\sim$500 steps) followed by slow refinement.}
  \label{fig:scale-shape}
\end{figure}

\begin{figure}[htbp]
  \centering
  \includegraphics[width=0.48\textwidth]{figures/baseline_class_geometry.png}
  \hfill
  \includegraphics[width=0.48\textwidth]{figures/baseline_topology.png}
  \caption{Baseline run. Left: class-geometry metrics (compactness,
  separation) over training. Right: topological summaries (H0/H1 entropy and
  lifetime sums, PMag) over training, computed on the normalized point
  cloud.}
  \label{fig:baseline-geom-topo}
\end{figure}

\subsection{Capacity sweep}

Test accuracy saturates quickly with latent dimension: $D{=}2$ lags at
$0.951$, but $D{=}4$ already reaches $0.971$ and $D \in \{8,16,32\}$ all sit
in a tight $0.975$--$0.976$ band (Table~\ref{tab:results},
Figure~\ref{fig:capacity}, top-left panel). By this metric alone, capacity
beyond $D{\approx}4$ is irrelevant.

The normalized geometry and topology tell a more nuanced story. The
normalized Fisher ratio is \emph{non-monotonic} in $D$: it is $5.56$ at
$D{=}2$, peaks at $6.49$ at $D{=}4$, and then \emph{decreases} through
$4.22$ ($D{=}8$), $3.62$ ($D{=}16$), to $3.55$ ($D{=}32$) -- a roughly
2$\times$ drop from its peak even as accuracy stays flat. H1 persistence
entropy and the H1 total-persistence sum ($L_1$) both increase monotonically
with $D$: entropy goes from $4.47$ ($D{=}2$) to $5.91$ ($D{=}32$), and
$L_1$ from $1.71$ to $12.33$, i.e.\ $D{=}2$ latents are topologically
impoverished (few, simple loop features) while higher-$D$ latents carry
substantially richer loop structure. PMag moves in the opposite direction,
decreasing from $11.46$ at $D{=}2$ to $7.93$ at $D{=}32$. The interpretation
consistent with all three signals: at low $D$, the encoder is forced to pack
class structure into a small number of dimensions, producing tight,
high-magnitude, topologically simple configurations with strong linear
(Fisher) separation once the obvious degrees of freedom are used up; at
higher $D$, the same classification job can be solved by spreading class
structure across more directions, which lowers the achievable Fisher ratio
per unit of within-class scatter while allowing richer loop topology to
develop (more independent 1-cycles with longer lifetimes) that linear
accuracy does not reward or penalize.

$D{=}16$ and $D{=}32$ are nearly indistinguishable on every shape metric
(fisher\_ratio $3.62$ vs.\ $3.55$; h1\_entropy $5.905$ vs.\ $5.912$;
h1\_lifetime\_sum $12.39$ vs.\ $12.33$; pmag $8.07$ vs.\ $7.93$), i.e.\
representational structure, not just accuracy, saturates by $D{=}16$:
doubling the latent dimension again to $D{=}32$ adds no further geometric or
topological structure that these metrics can detect.

\paragraph{Convergence timing: accuracy vs.\ geometry.} Accuracy and
geometry also converge on very different timescales. Defining
``converged'' as the first snapshot at which a metric reaches 90\% of its
own final-step value, test accuracy for $D\ge4$ converges almost
immediately -- within $4$--$6\%$ of the run's total steps (e.g.\ $D{=}16$:
step 100 of 2345, $4.3\%$) -- while the normalized Fisher ratio does not
reach 90\% of its final value until roughly $22$--$36\%$ of the way through
training (mean fraction of total steps: $35.5\%$ at $D{=}4$, $29.9\%$ at
$D{=}8$, $22.0\%$ at $D{=}16$, $27.7\%$ at $D{=}32$). In other words, the
linear readout stops improving almost as soon as training starts, but the
angular arrangement of class clusters keeps reorganizing for several times
longer. $D{=}2$ is the exception to the pattern: it is accuracy-bottlenecked
by its own capacity limit, so both accuracy and Fisher ratio converge late
and at similar times (mean fractions $13.5\%$ and $12.1\%$ respectively) --
there, unlike at higher $D$, the representation cannot finish separating
classes angularly without also continuing to push accuracy upward.

\begin{figure}[htbp]
  \centering
  \includegraphics[width=\textwidth]{figures/capacity_compare.png}
  \caption{Capacity sweep ($D \in \{2,4,8,16,32\}$, objective=probe, no
  normalization, mean $\pm$ std over 3 seeds). Test accuracy saturates by
  $D{\approx}4$, but fisher\_ratio is non-monotonic (peaking near $D{=}4$)
  and H1 entropy / H1 lifetime-sum increase monotonically with $D$ while
  PMag decreases with $D$ -- all metrics ripser/magnitude-based are computed
  on normalized latents.}
  \label{fig:capacity}
\end{figure}

\subsection{Objective contrast}

At fixed $D{=}16$, the training objective has a large effect on both
accuracy and geometry (Table~\ref{tab:results}, Figure~\ref{fig:objective}).
Readout accuracy is $0.976$ for probe and $0.977$ for probe\_recon, but only
$0.878$ for autoencoder -- confirming that reconstruction alone, with the
probe head detached from the encoder gradient, produces a representation
that is markedly less linearly separable by class even though the encoder
is free to use its full $D{=}16$ capacity.

Geometrically, supervision collapses the latent into clean, well-separated,
topologically simple class clusters: fisher\_ratio is $3.62$ (probe) $>
2.29$ (probe\_recon) $\gg 0.60$ (autoencoder), a roughly 6$\times$ spread
across objectives. The topology metrics invert this ordering: H1
lifetime-sum ($\alpha{=}1$) is $36.07$ for autoencoder, $18.26$ for
probe\_recon, and only $12.39$ for probe; H1 entropy follows the same
pattern ($6.55 > 6.17 > 5.91$). That is, reconstruction yields a diffuse,
topologically rich manifold -- rich loop structure not organized around
class identity -- while supervision prunes that topology away in favor of
compact, class-aligned clusters. PMag moves in \emph{opposite directions
over training} for the two pure objectives: it rises over training for
probe (reaching $8.07$) and falls for autoencoder (settling at $5.13$),
consistent with supervision concentrating mass into fewer, tighter
class-centered regions (higher effective ``resolution,'' larger magnitude at
a given scale) while reconstruction keeps the manifold spread out.
probe\_recon sits numerically between the two pure objectives on nearly
every metric, as expected for a loss that averages both signals.

\begin{figure}[htbp]
  \centering
  \includegraphics[width=\textwidth]{figures/objective_compare.png}
  \caption{Objective contrast ($D{=}16$, no normalization; probe vs.\
  autoencoder vs.\ probe\_recon, mean $\pm$ std over 3 seeds). Supervision
  (probe, probe\_recon) reaches high accuracy and high Fisher ratio with low
  topological complexity; pure reconstruction (autoencoder) trails in
  accuracy and Fisher ratio but develops much richer H1 topology.}
  \label{fig:objective}
\end{figure}

\subsection{Normalization contrast}
\label{sec:norm}

Constraining the latent with $\ell_2$ normalization (projection onto the
unit sphere) is essentially free in accuracy terms -- all three
normalization modes reach $0.976$--$0.977$ test accuracy
(Table~\ref{tab:results}, Figure~\ref{fig:norm}) -- but transforms the
geometry and topology. By construction, $\ell_2$ normalization drives
\texttt{raw\_mean\_norm} to exactly $1.0$ (vs.\ $23.6$ for no normalization
and $6.1$ for BatchNorm), removing the scale-inflation behavior described in
Section~\ref{sec:baseline} entirely. Normalized fisher\_ratio rises from
$3.62$ (none) to $4.71$ (BatchNorm) to $15.13$ ($\ell_2$) --
roughly a 4$\times$ jump over the unconstrained baseline -- while H1 entropy
drops from $5.91$ (none) to $5.74$ (BatchNorm) to $3.80$ ($\ell_2$), and
PMag jumps from $8.07$ (none) and $8.93$ (BatchNorm) to $52.73$ ($\ell_2$),
roughly a 6$\times$ increase. In short, forcing the latent onto the unit sphere produces much stronger
angular class separation and a much simpler, more concentrated
topology/magnitude profile, at zero accuracy cost.

It is important to be precise about what this gain is \emph{not}: the
Fisher ratio is mathematically invariant to isotropic rescaling of the
latent (Section~\ref{sec:methods}; verified empirically on every snapshot of
every run here, raw vs.\ normalized fisher\_ratio agree to $\sim10^{-15}$),
so the raw-scale inflation documented in Section~\ref{sec:baseline} cannot
by itself explain the unconstrained (``none'') run's lower Fisher ratio, and
symmetrically, $\ell_2$ normalization's $4\times$ Fisher-ratio gain over
``none'' cannot be a side effect of it also pinning \texttt{raw\_mean\_norm}
to exactly $1.0$. What $\ell_2$ and BatchNorm actually do is different: both
are applied \emph{during training}, inside the forward pass the encoder is
optimized against, and both are anisotropic, nonlinear, per-sample maps
(projection onto the sphere; per-feature standardization) rather than a
single global rescaling. Because gradients flow back through them, they
change what the encoder is incentivized to learn, not merely how its output
happens to be scaled after the fact. The large gains under $\ell_2$ --
Fisher ratio $3.6\to15.1$, compactness\_centroid $0.32\to0.11$, H1 entropy
$5.91\to3.80$, PMag $8.1\to52.7$ -- are therefore evidence of genuine
representational reshaping: the encoder learns to place classes at more
angularly separated locations on the sphere, in tighter per-class cones, at
zero accuracy cost. BatchNorm, despite also removing most of the raw-scale
inflation (\texttt{raw\_mean\_norm} $\to 6.1$), reshapes the representation
much less (Fisher ratio only $3.6\to4.7$), underscoring that it is the
specific geometry of the constraint (projection onto a sphere vs.\
per-feature standardization), not the removal of scale inflation per se,
that drives the $\ell_2$ result.

\begin{figure}[htbp]
  \centering
  \includegraphics[width=\textwidth]{figures/norm_compare.png}
  \caption{Normalization contrast ($D{=}16$, objective=probe; none vs.\
  $\ell_2$ vs.\ BatchNorm, mean $\pm$ std over 3 seeds). All three reach
  essentially the same test accuracy, but $\ell_2$ normalization produces a
  far higher Fisher ratio, lower H1 entropy, and much larger PMag than
  either alternative.}
  \label{fig:norm}
\end{figure}

\subsection{Latent-cloud evolution: trajectory grids}

Figure~\ref{fig:trajectories} contrasts the evolution of the (PCA,
fixed-basis) latent point cloud across four representative configurations:
low-capacity supervised ($D{=}2$, probe), the $D{=}16$ probe baseline, the
$D{=}16$ autoencoder, and the $D{=}16$ $\ell_2$-normalized probe. The
$D{=}2$ cloud is visibly constrained to occupy the full 2D plane with
classes packed into adjacent wedges; the $D{=}16$ probe baseline (projected
onto its leading PCs) shows tighter, more separated blobs; the autoencoder
cloud remains visibly more diffuse and overlapping throughout training,
consistent with its low Fisher ratio and high topological complexity; and
the $\ell_2$-normalized cloud shows classes resolving into angularly
separated clusters on the projected unit circle/sphere, consistent with its
much higher Fisher ratio.

\begin{figure}[htbp]
  \centering
  \begin{subfigure}[t]{0.48\textwidth}
    \includegraphics[width=\textwidth]{figures/traj_probe_d2.png}
    \caption{probe, $D{=}2$}
  \end{subfigure}
  \hfill
  \begin{subfigure}[t]{0.48\textwidth}
    \includegraphics[width=\textwidth]{figures/traj_probe_d16.png}
    \caption{probe, $D{=}16$}
  \end{subfigure}
  \par\bigskip
  \begin{subfigure}[t]{0.48\textwidth}
    \includegraphics[width=\textwidth]{figures/traj_autoencoder_d16.png}
    \caption{autoencoder, $D{=}16$}
  \end{subfigure}
  \hfill
  \begin{subfigure}[t]{0.48\textwidth}
    \includegraphics[width=\textwidth]{figures/traj_l2_d16.png}
    \caption{probe, $D{=}16$, $\ell_2$-normalized}
  \end{subfigure}
  \caption{Trajectory grids (PCA, fixed basis) showing the evolution of the
  probe-set latent point cloud over training (colored by class) for four
  representative runs. Low-capacity ($D{=}2$) and reconstruction-only
  (autoencoder) latents remain visibly more entangled across classes than
  the $D{=}16$ supervised baseline; $\ell_2$ normalization produces the
  most visually separated class structure.}
  \label{fig:trajectories}
\end{figure}

\section{Discussion}

Taken together, these results show that final-step test accuracy is a very
incomplete summary of what a classifier's latent space is doing. Three
instances of this stand out:

\begin{itemize}
  \item \textbf{Scale masks shape.} Raw distance-based statistics in the
  baseline run are dominated by an isotropic $\sim$40$\times$ norm inflation
  that has nothing to do with how class structure reorganizes; only after
  normalizing out this scale does the real reorganization -- a sharp
  elbow within the first epoch, then slow refinement -- become visible.
  Any study of representation geometry over training that does not control
  for this will conflate the two.

  \item \textbf{Capacity beyond the point needed for accuracy still
  matters.} Test accuracy is flat for $D \ge 4$, but the normalized Fisher
  ratio, H1 entropy, H1 lifetime sum, and PMag all keep changing
  substantially across $D \in \{4,8,16,32\}$. Low-dimensional latents are
  pushed into tight, high-Fisher-ratio, topologically simple, high-magnitude
  configurations; higher-dimensional latents spread the same class
  information across more directions, at the cost of Fisher ratio but with
  richer loop topology. A linear probe cannot see this difference, but it is
  a real difference in how the representation is organized, and could
  matter for transfer, robustness, or downstream nonlinear use of the
  representation.

  \item \textbf{Accuracy converges early; geometry keeps reorganizing.} In
  the capacity sweep, test accuracy for $D\ge4$ reaches 90\% of its
  final value within $\sim$4--6\% of training steps, while the (scale-
  invariant) Fisher ratio does not reach 90\% of its final value until
  $\sim$22--36\% of the way through training. The linear readout stops
  improving almost immediately; the angular arrangement of class clusters
  continues to reorganize for several times longer. Any study that stops
  measuring representation quality once accuracy plateaus would miss most
  of this reorganization.

  \item \textbf{Normalization reshapes representations, it does not just
  rescale them.} The normalization contrast is the clearest evidence that
  accuracy alone under-describes representation quality:
  $\ell_2$-normalizing the latent costs nothing in accuracy but roughly
  quadruples the (scale-invariant) Fisher ratio and increases PMag by
  $\sim$6$\times$ relative to the unconstrained baseline. Because the Fisher
  ratio is provably invariant to isotropic rescaling, this is not a
  byproduct of $\ell_2$ normalization incidentally pinning the raw norm to
  $1$; it reflects a real change in how the encoder arranges classes
  angularly, driven by training with a nonlinear, anisotropic constraint in
  the forward pass. BatchNorm, which also removes most of the raw-scale
  inflation but reshapes the representation far less (Fisher ratio
  $3.6\to4.7$ vs.\ $3.6\to15.1$ for $\ell_2$), makes the same point from the
  other direction: removing scale inflation is not sufficient on its own to
  produce better-separated class structure, so the specific (spherical)
  geometry of the $\ell_2$ constraint is doing the work.
\end{itemize}

\paragraph{Limitations.} This study uses MNIST only, with a small MLP
encoder (784-256-64-$D$) and few training epochs (5); conclusions about
capacity, objective, and normalization effects may not transfer to deeper
architectures, harder datasets, or longer training schedules. Persistent
homology is computed with \texttt{ripser} on point clouds of $N{=}1000$
probe points per snapshot, which is adequate for qualitative H0/H1 summaries
at this scale but becomes increasingly expensive (and the resulting diagrams
noisier) for larger $N$ or higher homological degree; we restrict to H0/H1
throughout. Each configuration is evaluated over 3 random seeds, which is
enough to see consistent trends and non-overlapping error bands in the
figures above but is a modest sample for statistical claims beyond ``clearly
separated vs.\ not.'' Finally, all topological and magnitude-function
metrics are computed on mean-pairwise-distance-normalized latents to isolate
shape from scale; this is one principled normalization choice among several
(e.g.\ per-dimension z-scoring is also implemented and is anisotropic, hence
not shape-preserving) and the quantitative values reported here are specific
to it.

\section{Results Table}

\begin{sloppypar}
Table~\ref{tab:results} reports per-configuration final-step metrics, mean
$\pm$ std over 3 seeds, aggregated from the 27-run sweep described in
Section~1. All geometric/topological columns (Fisher ratio, H1 entropy, H1
lifetime sum, PMag) are computed on normalized latents; raw\_mean\_norm is
the un-normalized scale.
\end{sloppypar}

\begin{table}[htbp]
\centering
\resizebox{\textwidth}{!}{%
\begin{tabular}{llrrrrrrr}
\toprule
Objective & Norm & $D$ & test\_acc & fisher\_ratio & h1\_entropy & h1\_lifesum$_{\alpha=1}$ & pmag & raw\_mean\_norm \\
\midrule
probe        & none & 2  & $0.951 \pm 0.001$ & $5.56 \pm 0.23$  & $4.47 \pm 0.06$ & $1.71 \pm 0.16$  & $11.46 \pm 0.24$ & $53.75 \pm 2.45$ \\
probe        & none & 4  & $0.971 \pm 0.000$ & $6.49 \pm 0.45$  & $5.49 \pm 0.04$ & $6.60 \pm 0.54$  & $10.21 \pm 0.30$ & $23.27 \pm 2.21$ \\
probe        & none & 8  & $0.975 \pm 0.002$ & $4.22 \pm 0.31$  & $5.84 \pm 0.02$ & $11.53 \pm 0.75$ & $8.13 \pm 0.32$  & $23.61 \pm 1.56$ \\
probe        & none & 16 & $0.976 \pm 0.001$ & $3.62 \pm 0.12$  & $5.91 \pm 0.04$ & $12.39 \pm 0.28$ & $8.07 \pm 0.16$  & $23.58 \pm 0.44$ \\
probe        & none & 32 & $0.975 \pm 0.002$ & $3.55 \pm 0.05$  & $5.91 \pm 0.01$ & $12.33 \pm 0.09$ & $7.93 \pm 0.09$  & $24.91 \pm 0.37$ \\
\midrule
autoencoder  & none & 16 & $0.878 \pm 0.003$ & $0.60 \pm 0.02$  & $6.55 \pm 0.02$ & $36.07 \pm 1.13$ & $5.13 \pm 0.05$  & $22.50 \pm 1.34$ \\
probe\_recon & none & 16 & $0.977 \pm 0.000$ & $2.29 \pm 0.07$  & $6.17 \pm 0.02$ & $18.26 \pm 0.78$ & $6.50 \pm 0.12$  & $27.51 \pm 0.73$ \\
\midrule
probe        & $\ell_2$ & 16 & $0.976 \pm 0.002$ & $15.13 \pm 0.81$ & $3.80 \pm 0.12$ & $3.52 \pm 0.60$  & $52.73 \pm 1.61$ & $1.00 \pm 0.00$ \\
probe        & BN       & 16 & $0.977 \pm 0.001$ & $4.71 \pm 0.27$  & $5.74 \pm 0.04$ & $9.64 \pm 0.24$  & $8.93 \pm 0.07$  & $6.12 \pm 0.13$ \\
\bottomrule
\end{tabular}%
}
\caption{Final-step metrics, mean $\pm$ std over 3 seeds, for all 9
configurations in the 27-run sweep (capacity sweep: top block; objective
contrast: middle block, sharing its $D{=}16$/none/probe row with the top
block; normalization contrast: bottom block, sharing the same row).
Source: \texttt{report\_assets/summary\_all.csv}.}
\label{tab:results}
\end{table}

\section*{References}
\begin{itemize}
  \item[] Leinster, T. (2013). The magnitude of metric spaces. \emph{Documenta Mathematica}, 18, 857--905.
  \item[] Govc, D. \& Hepworth, R. (2021). Persistent magnitude. \emph{Journal of Pure and Applied Algebra}, 225(3), 106517. arXiv:2010.09776.
  \item[] Rucco, M. et al. Characterisation of the idling behaviour of neural networks via persistence entropy.
  \item[] Atienza, N., Gonzalez-Diaz, R., \& Soriano-Trigueros, M. (2020). On the stability of persistent entropy and new summary functions for TDA.
\end{itemize}

\end{document}
